% GENERATED FILE. Edit canonical lessons, then run npm run generate:books.
% canonical-source-hash: fnv1a64:74cd356fbb240092
% canonical-lessons: TE-C424-ottidi, TE-C424-balavantam, TE-C424-abhyarthincu, TE-C424-edirincu, TE-C424-upasamharincu

\chapter{Pressure and Resistance}
\label{ch:te-pressure-and-resistance}
\hlchaptermodality{424}

\begin{chapteropening}
\textbf{By the end of this chapter:} \emph{I can talk about pressure, force, appeals and resistance.}

Complete the last lesson of chapter 424: I can talk about pressure, force, appeals and resistance.
\end{chapteropening}

\section[\texorpdfstring{\te{ఒత్తిడి} (ottiḍi)}{ఒత్తిడి (ottiḍi)}]{\te{ఒత్తిడి} (ottiḍi) — pressure; strain}

\label{lesson:TE-C424-ottidi}



\begin{quote}
Before the new one: say the Telugu for to be enough, then the Telugu for to make use of.
\end{quote}

\subsection*{You'll want to know: \te{ఒత్తిడి}}
\textbf{\te{ఒత్తిడి}} — \emph{ottiḍi} — \textquotedblleft{}pressure; strain\textquotedblright{}.

\begin{quote}
\te{పరీక్షల} \te{ముందు} \te{విద్యార్థుల} \te{మీద} \te{చాలా} \te{ఒత్తిడి} \te{ఉంటుంది}.
\end{quote}

\begin{quote}
\emph{parīkṣala mundu vidyārthula mīda cālā ottiḍi uṇṭundi.}
\end{quote}

\begin{quote}
Before the exams students are under a lot of pressure.
\end{quote}

\te{ఒత్తిడి} \te{తెచ్చు}, to put pressure on someone.

\begin{cousinweb}
\textbf{\te{ఒత్తిడి}} (\emph{ottiḍi}) is native Dravidian, built on \emph{ottu}, \textquotedblleft{}to press, to squeeze; pressure\textquotedblright{} (DEDR 1021).
\end{cousinweb}

\begin{grammarlens}[title={What you've built}]
One more word for saying what happens if.
\end{grammarlens}

\subsection*{Your turn}
\begin{itemize}
\raggedright
  \item \emph{Say it:} \emph{ottiḍi}
  \item \emph{Say it:} \emph{ottiḍi}, once more
  \item \emph{Say it:} \emph{parīkṣala mundu vidyārthula mīda cālā ottiḍi uṇṭundi.} — Before the exams students are under a lot of pressure.
  \item \emph{Recall:} say the Telugu for to be enough, then the Telugu for to make use of, then say \emph{ottiḍi} again
\end{itemize}

\begin{tcolorbox}[breakable,colback=teal!4,colframe=teal!35!black,title={Before you move on}]
What does \te{ఒత్తిడి} mean? (\textquotedblleft{}Pressure; strain\textquotedblright{}.) Say it once more.
\end{tcolorbox}
\section[\texorpdfstring{\te{బలవంతం} (balavantaṁ)}{బలవంతం (balavantaṁ)}]{\te{బలవంతం} (balavantaṁ) — force; compulsion}

\label{lesson:TE-C424-balavantam}



\begin{quote}
Before the new one: say the Telugu for to make use of, then the Telugu for pressure.
\end{quote}

\subsection*{You'll want to know: \te{బలవంతం}}
\textbf{\te{బలవంతం}} — \emph{balavantaṁ} — \textquotedblleft{}force; compulsion\textquotedblright{}.

\begin{quote}
\te{ఎవరినీ} \te{బలవంతంగా} \te{పనిలో} \te{పెట్టకూడదు}.
\end{quote}

\begin{quote}
\emph{evarinī balavantaṅgā panilō peṭṭakūḍadu.}
\end{quote}

\begin{quote}
No one should be forced into work.
\end{quote}

\te{బలవంతంగా}, by force; \te{బలవంతం} \te{చేయు}, to force.

\begin{cousinweb}
\textbf{\te{బలవంతం}} (\emph{balavantaṁ}) is built on Sanskrit \emph{balavant}, \textquotedblleft{}strong\textquotedblright{} (CDIAL 9166), from \emph{bala}, \textquotedblleft{}strength\textquotedblright{}, the word in \te{బలం}.
\end{cousinweb}

\begin{grammarlens}[title={What you've built}]
One more word for saying what happens if.
\end{grammarlens}

\subsection*{Your turn}
\begin{itemize}
\raggedright
  \item \emph{Say it:} \emph{balavantaṁ}
  \item \emph{Say it:} \emph{balavantaṁ}, once more
  \item \emph{Say it:} \emph{evarinī balavantaṅgā panilō peṭṭakūḍadu.} — No one should be forced into work.
  \item \emph{Recall:} say the Telugu for to make use of, then the Telugu for pressure, then say \emph{balavantaṁ} again
\end{itemize}

\begin{tcolorbox}[breakable,colback=teal!4,colframe=teal!35!black,title={Before you move on}]
What does \te{బలవంతం} mean? (\textquotedblleft{}Force; compulsion\textquotedblright{}.) Say it once more.
\end{tcolorbox}
\section[\texorpdfstring{\te{అభ్యర్థించు} (abhyarthiñcu)}{అభ్యర్థించు (abhyarthiñcu)}]{\te{అభ్యర్థించు} (abhyarthiñcu) — to appeal; to request formally}

\label{lesson:TE-C424-abhyarthincu}



\begin{quote}
Before the new one: say the Telugu for pressure, then the Telugu for force.
\end{quote}

\subsection*{You'll want to know: \te{అభ్యర్థించు}}
\textbf{\te{అభ్యర్థించు}} — \emph{abhyarthiñcu} — \textquotedblleft{}to appeal; to request formally\textquotedblright{}.

\begin{quote}
\te{వాళ్ళు} \te{సహాయం} \te{కోసం} \te{ప్రభుత్వాన్ని} \te{అభ్యర్థించారు}.
\end{quote}

\begin{quote}
\emph{vāḷḷu sahāyaṁ kōsaṁ prabhutvānni abhyarthiñcāru.}
\end{quote}

\begin{quote}
They appealed to the government for help.
\end{quote}

Past (I …): \te{అభ్యర్థించాను}, \emph{abhyarthiñcānu}.

The noun is \te{అభ్యర్థన}, a request; \te{అభ్యర్థి}, a candidate, is the same family.

\begin{cousinweb}
\textbf{\te{అభ్యర్థించు}} (\emph{abhyarthiñcu}) is built on Sanskrit \emph{abhyarthana}, \textquotedblleft{}asking, requesting\textquotedblright{} (Monier-Williams), with the Telugu verb ending -\te{ఇంచు}.
\end{cousinweb}

\begin{grammarlens}[title={What you've built}]
One more word for saying what happens if.
\end{grammarlens}

\subsection*{Your turn}
\begin{itemize}
\raggedright
  \item \emph{Say it:} \emph{abhyarthiñcu}
  \item \emph{Say it:} \emph{abhyarthiñcu}, once more
  \item \emph{Say it:} \emph{vāḷḷu sahāyaṁ kōsaṁ prabhutvānni abhyarthiñcāru.} — They appealed to the government for help.
  \item \emph{Recall:} say the Telugu for pressure, then the Telugu for force, then say \emph{abhyarthiñcu} again
\end{itemize}

\begin{tcolorbox}[breakable,colback=teal!4,colframe=teal!35!black,title={Before you move on}]
What does \te{అభ్యర్థించు} mean? (\textquotedblleft{}To appeal; to request formally\textquotedblright{}.) Say it once more.
\end{tcolorbox}
\section[\texorpdfstring{\te{ఎదిరించు} (ediriñcu)}{ఎదిరించు (ediriñcu)}]{\te{ఎదిరించు} (ediriñcu) — to stand up to; to resist}

\label{lesson:TE-C424-edirincu}



\begin{quote}
Before the new one: say the Telugu for force, then the Telugu for to appeal.
\end{quote}

\subsection*{You'll want to know: \te{ఎదిరించు}}
\textbf{\te{ఎదిరించు}} — \emph{ediriñcu} — \textquotedblleft{}to stand up to; to resist\textquotedblright{}.

\begin{quote}
\te{ఆమె} \te{అన్యాయాన్ని} \te{ధైర్యంగా} \te{ఎదిరించింది}.
\end{quote}

\begin{quote}
\emph{āme anyāyānni dhairyaṅgā ediriñcindi.}
\end{quote}

\begin{quote}
She bravely stood up to injustice.
\end{quote}

Past (I …): \te{ఎదిరించాను}, \emph{ediriñcānu}.

\te{ఎదురు}, opposite, in front, is the same root.

\begin{cousinweb}
\textbf{\te{ఎదిరించు}} (\emph{ediriñcu}) is native Dravidian: DEDR 795 lists Telugu \emph{edirincu}, \textquotedblleft{}to oppose, to resist\textquotedblright{}, beside \emph{eduru}, \textquotedblleft{}the front, the opposite\textquotedblright{}.
\end{cousinweb}

\begin{grammarlens}[title={What you've built}]
One more word for saying what happens if.
\end{grammarlens}

\subsection*{Your turn}
\begin{itemize}
\raggedright
  \item \emph{Say it:} \emph{ediriñcu}
  \item \emph{Say it:} \emph{ediriñcu}, once more
  \item \emph{Say it:} \emph{āme anyāyānni dhairyaṅgā ediriñcindi.} — She bravely stood up to injustice.
  \item \emph{Recall:} say the Telugu for force, then the Telugu for to appeal, then say \emph{ediriñcu} again
\end{itemize}

\begin{tcolorbox}[breakable,colback=teal!4,colframe=teal!35!black,title={Before you move on}]
What does \te{ఎదిరించు} mean? (\textquotedblleft{}To stand up to; to resist\textquotedblright{}.) Say it once more.
\end{tcolorbox}
\section[\texorpdfstring{\te{ఉపసంహరించు} (upasaṁhariñcu)}{ఉపసంహరించు (upasaṁhariñcu)}]{\te{ఉపసంహరించు} (upasaṁhariñcu) — to withdraw; to take back}

\label{lesson:TE-C424-upasamharincu}



\begin{quote}
Before the new one: say the Telugu for to appeal, then the Telugu for to stand up to.
\end{quote}

\subsection*{You'll want to know: \te{ఉపసంహరించు}}
\textbf{\te{ఉపసంహరించు}} — \emph{upasaṁhariñcu} — \textquotedblleft{}to withdraw; to take back\textquotedblright{}.

\begin{quote}
\te{అతను} \te{తన} \te{దరఖాస్తును} \te{ఉపసంహరించాడు}.
\end{quote}

\begin{quote}
\emph{atanu tana darakhāstunu upasaṁhariñcāḍu.}
\end{quote}

\begin{quote}
He withdrew his application.
\end{quote}

Past (I …): \te{ఉపసంహరించాను}, \emph{upasaṁhariñcānu}.

The noun is \te{ఉపసంహరణ}, a withdrawal.

\begin{cousinweb}
\textbf{\te{ఉపసంహరించు}} (\emph{upasaṁhariñcu}) is built on Sanskrit \emph{upasaṁhāra}, \textquotedblleft{}the act of withdrawing, withholding\textquotedblright{} (Monier-Williams), with the Telugu verb ending -\te{ఇంచు}.
\end{cousinweb}

\begin{grammarlens}[title={What you've built}]
One more word for saying what happens if.
\end{grammarlens}

\subsection*{Your turn}
\begin{itemize}
\raggedright
  \item \emph{Say it:} \emph{upasaṁhariñcu}
  \item \emph{Say it:} \emph{upasaṁhariñcu}, once more
  \item \emph{Say it:} \emph{atanu tana darakhāstunu upasaṁhariñcāḍu.} — He withdrew his application.
  \item \emph{Recall:} say the Telugu for to appeal, then the Telugu for to stand up to, then say \emph{upasaṁhariñcu} again
\end{itemize}

\begin{tcolorbox}[breakable,colback=teal!4,colframe=teal!35!black,title={Before you move on}]
What does \te{ఉపసంహరించు} mean? (\textquotedblleft{}To withdraw; to take back\textquotedblright{}.) Say it once more.
\end{tcolorbox}
